\section{Theoretical Foundations of the Arithmetic Mismatch}
\label{sec:theoretical}

This section sets up the formal cost framework against which the rest of the thesis is measured. We name the dimensions along which the static-circuit and zkVM regimes will be compared (Section~\ref{ssec:cost-dimensions}), formalise the \emph{field-mismatch tax} as the quantitative manifestation of the arithmetic mismatch (Section~\ref{ssec:field-mismatch-tax}), formalise \emph{update churn} as a deployment-lifecycle dimension (Section~\ref{ssec:update-churn}), and decompose the cost of the BDEC CreGen and ShowCre relations into the components that dominate either regime (Section~\ref{ssec:cost-decomposition}).

\subsection{Cost Dimensions}\label{ssec:cost-dimensions}

We measure each system regime along four dimensions.

\begin{description}
  \item[Time.] Prover and verifier wall-clock time, measured in seconds on the target hardware. For the prover this is the time from invocation of $\Pi.\Prove$ (or $\Pi_{\mathsf{vm}}.\Prove$) until the receipt is produced; for the verifier it is the time from invocation of the verification predicate until the binary outcome is returned.
  \item[Memory.] Peak resident-set size (peak RSS) of the prover process on the target hardware, measured in gigabytes. Verifier memory is reported when it materially differs from prover memory; otherwise it is dominated by the receipt size.
  \item[Proof size.] The size of the produced receipt, measured in kilobytes for the relevant receipt class. We report this separately for the succinct STARK receipt and the SNARK-wrapped receipt, the two receipt regimes of Section~\ref{sec:prelim}, where each is produced.
  \item[Update churn.] The footprint of artefacts that must be regenerated and re-audited when the verification predicate changes by one localised modification. We formalise this dimension in Section~\ref{ssec:update-churn} below.
\end{description}

The first three dimensions are standard in the zkSNARK and zkVM literature~\cite{cryptoeprint:2024/868,10.1007/978-981-95-2961-2_6,bruestle2023risczero,succinct2024sp1,cryptoeprint:2023/1217}. The fourth dimension is introduced in this thesis as a deployment-lifecycle accounting frame. Unlike the first three it is a footprint characterisation rather than a performance metric, and Section~\ref{sec:evaluation} shows that it does not translate into a wall-clock advantage against the non-preprocessing static baseline.

\subsection{The Field-Mismatch Tax}\label{ssec:field-mismatch-tax}

The verification predicate of a post-quantum signature such as Loquat or PLUM is dominated by algebraic-hash (Griffin) permutations together with arithmetic over a prime field $\mathbb{F}_p$ of large characteristic (in PLUM, a nominal $192$-bit / concrete $199$-bit prime; see Section~\ref{sec:prelim}). Concretely, PLUM~\cite[\S4.2]{10.1007/978-981-95-2961-2_6} reports that hash evaluations account for $105{,}952$ of the $116{,}285$ R1CS constraints ($\approx 91\%$) of a PLUM-128 verification, the field-arithmetic constraints making up the remaining $\approx 9\%$; the tax below is carried by this hash-dominated, large-field workload as a whole, and not by field arithmetic on its own. This $91\%$ is a constraint count in a relation-specific zkSNARK and is not a measured zkVM-cycle share: it does not, on its own, fix an Amdahl ceiling on the zkVM cost, and the measured zkVM-cycle attribution is reported separately in Section~\ref{sec:evaluation}. When such a predicate is proved by a relation-specific zkSNARK over $\mathbb{F}_p$, every field operation in the predicate is a single constraint over the same field. When it is proved by a zkVM whose prover field $\mathbb{F}_{p_{\mathsf{vm}}}$ is small (e.g.\ BabyBear, $p_{\mathsf{vm}} = 2^{31} - 2^{27} + 1$), each $\mathbb{F}_p$ operation must be emulated by software running on the zkVM's RISC-V instruction set, which in turn is encoded in $\mathbb{F}_{p_{\mathsf{vm}}}$. The number of native RISC-V instructions required to emulate one $\mathbb{F}_p$ operation by multi-limb arithmetic is a fixed function of the bit-widths of $p$ and the limb size of the underlying integer encoding.

The emulation tax this definition names is characterised qualitatively in the recent zkVM systematisation~\cite{yang2026sokzkvm}. The definition below is a cost-accounting frame that names, in one symbol, the per-operation overhead the precompile suite is built to remove.
\begin{definition}[Field-mismatch tax]\label{def:fmt}
Let $\mathsf{op}$ be a field operation over $\mathbb{F}_p$ used inside a verification predicate $\Pred$, and let $\mathsf{cyc}_{\mathsf{em}}(\mathsf{op})$ denote the number of zkVM guest cycles required to emulate $\mathsf{op}$ by software on a zkVM whose native field is $\mathbb{F}_{p_{\mathsf{vm}}}$. Let $\mathsf{cyc}_{\mathsf{nat}}(\mathsf{op})$ denote the number of zkVM guest cycles that the same operation would consume if the prover field were $\mathbb{F}_p$ itself, so that no multi-limb emulation is needed; under this idealisation $\mathsf{cyc}_{\mathsf{nat}}(\mathsf{op})$ is the cycle count of a single native field instruction. Because the target zkVMs expose no native $\mathbb{F}_p$ instruction, $\mathsf{cyc}_{\mathsf{nat}}$ is a conceptual normaliser, not a directly measurable quantity; the empirically measured quantity is the emulated cost $\mathsf{Cyc}_{\mathsf{em}}$ reported in Section~\ref{sec:evaluation}. Both quantities are in the same unit, so their ratio is dimensionless. The \emph{field-mismatch tax} of $\mathsf{op}$ is
\[
  \mathsf{FMT}(\mathsf{op})
  =
  \frac{\mathsf{cyc}_{\mathsf{em}}(\mathsf{op})}{\mathsf{cyc}_{\mathsf{nat}}(\mathsf{op})}.
\]
Writing $\mathsf{ops}(\Pred)$ for the multiset of field operations executed by an honest evaluation of $\Pred$, the cumulative \emph{emulated cost} of $\Pred$ is the distinct quantity
\[
  \mathsf{Cyc}_{\mathsf{em}}(\Pred)
  =
  \sum_{\mathsf{op}\in\mathsf{ops}(\Pred)}
  \mathsf{cyc}_{\mathsf{em}}(\mathsf{op}).
\]
\end{definition}

A precompile, as defined in Section~\ref{sec:prelim}, replaces $\mathsf{cyc}_{\mathsf{em}}(\mathsf{op})$ for a designated operation by a single zkVM syscall, whose cost from the guest's perspective is one cycle plus the constant overhead of the syscall ABI. Writing $\mathsf{Pre}\subseteq\mathsf{ops}(\Pred)$ for the set of operations covered by precompiles and $\mathsf{cyc}_{\mathsf{syscall}}$ for the per-syscall overhead, the cumulative emulated cost of $\Pred$ under a precompile-enabled zkVM is therefore
\[
  \mathsf{Cyc}_{\mathsf{em}}^{\mathsf{pre}}(\Pred)
  =
  \sum_{\mathsf{op}\in\mathsf{ops}(\Pred)\setminus\mathsf{Pre}}
  \mathsf{cyc}_{\mathsf{em}}(\mathsf{op})
  + |\{\mathsf{op}\in\mathsf{ops}(\Pred)\cap\mathsf{Pre}\}|
  \cdot \mathsf{cyc}_{\mathsf{syscall}}.
\] The performance benefit of adding a primitive to the precompile set is, by this accounting, the difference $\mathsf{cyc}_{\mathsf{em}}(\mathsf{op}) - \mathsf{cyc}_{\mathsf{syscall}}$ summed over its occurrences. This is a \emph{guest-trace} reduction: it counts cycles removed from the emulated execution, not the prover's wall-clock cost. The precompile shifts those cycles into a relation-specific AIR whose proving cost is charged on the host side, so a positive guest-cycle benefit need not translate into a faster prover, and Section~\ref{sec:evaluation} reports a primitive (Griffin) for which it does not.

\subsection{Update Churn}\label{ssec:update-churn}

\begin{definition}[Localised verification-predicate change]\label{def:lpc}
A \emph{localised verification-predicate change} is a modification to the verification predicate $\Pred$ that can be described as one of the following: substitution of one signature scheme for another at matching security level; retuning of a single numerical security parameter (e.g.\ $\lambda$, $t$, $\eta$); a change to the presentation \emph{structure} of a showing relation (the number of credentials shown, or the set of disclosed attributes); or a one-line change to the encoding of a public input. A change that only alters which condition the verifier checks on the \emph{disclosed} attributes, a presentation \emph{policy} at fixed structure, is \emph{not} a verification-predicate change: following base BDEC it is evaluated outside the proof relation (Section~\ref{sec:prelim}) and touches neither $\Pred$ nor its compiled artefacts.
\end{definition}

The deployment-lifecycle cost of a localised change is dominated by the artefacts that must be regenerated and re-audited as a consequence, rather than by the change itself. We formalise this cost as follows.

\begin{definition}[Update churn]\label{def:update-churn}
Fix a deployment $\mathcal D$ that proves $\Pred$ in either the static-circuit or the zkVM regime, and let $\Delta$ denote a localised verification-predicate change. The \emph{update-churn footprint} of $\Delta$ in $\mathcal D$ is the tuple
\[
  \mathsf{Churn}(\Delta,\mathcal D)
  =
  \bigl(
    N_{\mathsf{src}}(\Delta,\mathcal D),
    N_{\mathsf{circ}}(\Delta,\mathcal D),
    N_{\mathsf{param}}(\Delta,\mathcal D),
    N_{\mathsf{audit}}(\Delta,\mathcal D)
  \bigr),
\]
where
\begin{description}
  \item[$N_{\mathsf{src}}$] is the number of source-code files that must be modified in response to $\Delta$.
  \item[$N_{\mathsf{circ}}$] is the number of arithmetic circuits or AIR programs that must be regenerated, counted as the distinct circuit or AIR artefacts whose serialised form changes between the pre- and post-$\Delta$ builds.
  \item[$N_{\mathsf{param}}$] is the number of proof-system parameters (e.g.\ commitment parameters, R1CS-shape-dependent constants) that must be re-derived.
  \item[$N_{\mathsf{audit}}$] is the number of distinct components requiring re-audit before redeployment, counted at the granularity of the deployment's existing security-review boundary. This is an audit-\emph{footprint} count, not a weighted effort model: it reports how many components re-enter the review boundary, not the person-time or risk class of reviewing them.
\end{description}
\end{definition}

For the static-circuit regime, a localised change to $\Pred$ typically yields a non-trivial $N_{\mathsf{circ}}$ (the R1CS must be recompiled), a non-trivial $N_{\mathsf{param}}$ (constants that depend on the constraint count must be re-derived), and a non-trivial $N_{\mathsf{audit}}$ (the new circuit shape must be audited).
For the zkVM regime, the same change always yields a non-trivial $N_{\mathsf{src}}$ (the guest program is edited) and zero regenerations of the \emph{universal} relation $R_{\mathsf{vm}}$, which is unchanged. The remaining components depend on whether the change touches a precompiled primitive. The bespoke $\mathsf{Griffin}_{\mathbb{F}_p^{192}}$ AIR is itself a relation-specific circuit; a change that alters the hash family or the field width therefore regenerates and re-audits that AIR, so $N_{\mathsf{circ}}\ge 1$, $N_{\mathsf{param}}\ge 1$, and $N_{\mathsf{audit}}\ge 1$. The zero-circuit, zero-parameter case applies only to changes that leave every precompiled primitive intact (e.g.\ raising the presentation count $k$ from $1$ to $2$), for which $N_{\mathsf{circ}}=0$, $N_{\mathsf{param}}=0$, and $N_{\mathsf{audit}}$ is confined to the changed guest source. The configuration that would need zero new circuits for \emph{any} change is the pure-emulation zkVM without precompiles, which Section~\ref{sec:evaluation} measures as infeasible (does not finish within the resource bound on the target hardware).

The empirical evaluation of Section~\ref{sec:evaluation} reports each of these four components for both regimes under a fixed set of localised changes that we exhibit as benchmarks: signature substitution (Loquat $\to$ PLUM), security-parameter retuning ($\lambda=80\to128$), and a presentation-structure change (raising the credential count $k$ from $1$ to $2$).

\subsection{Cost Decomposition of BDEC Relations}\label{ssec:cost-decomposition}

We now apply the cost framework above to the two BDEC relations of interest. We treat each relation as a multiset of cost-bearing components and compute the cumulative cost under each regime.

\paragraph{CreGen.}
The BDEC CreGen relation $R_{\mathsf{cre}}^{\Sig}$ defined in Section~\ref{sec:prelim} is the base two-verification relation: two signature-verification predicates $\Pred_{\Verify}^{\Sig,\ppSig}$ under a hidden user public key. The credential digest $\rho_{U,TA}=H(\cdot)$ is computed at issuance, and ledger membership is checked outside the proof, as in ShowCre; the full acceptance condition is therefore the conjunction of the in-proof signature-verification predicates with the verifier-side membership check performed by $\BDEC.\mathsf{CreVer}$/$\BDEC.\mathsf{ShowVer}$ (Section~\ref{sec:prelim}), so a proof for valid signatures whose digest is absent from the accepted ledger is rejected outside the proof.
Of these, the two $\Pred_{\Verify}^{\Sig,\ppSig}$ predicates dominate the cost. Inside PLUM verification (Section~\ref{sec:prelim}), the load-bearing components are (i) the $t$-th power-residue PRF symbol checks, (ii) the BCS Merkle-commitment recomputations, and (iii) the low-degree-test (LDT) sub-protocol invocations.
The cost of (ii) and (iii) is dominated by the algebraic hash function $H_{\mathsf{Griffin}}$ over $\mathbb{F}_p$, which is precisely the primitive absorbed by the $\mathsf{Griffin}_{\mathbb{F}_p^{192}}$ precompile.

\paragraph{ShowCre.}
The ShowCre proof circuit contains $k+2$ signature-verification predicates (where $k$ is the number of credentials shown): $k$ pseudonym-ownership checks, one verifier-facing pseudonym check, and one shown-credential check (Section~\ref{sec:prelim}). Ledger membership, non-revocation, and the disclosure predicate are checked by the relying party outside the proof.
For modest $k$ (the design target is $k\le 4$), the cost is dominated by these signature-verification predicates and the algebraic-hash operations inside their Merkle and low-degree-test components.

In both relations, the field-mismatch tax of Definition~\ref{def:fmt} is concentrated in the signature-verification predicates: every Griffin permutation involved in a Merkle path or in PLUM's internal hashing is a tax-bearing $\mathbb{F}_p^{192}$ operation under a zkVM with a small native field.

The decomposition locates the dominant term of both relations in the Griffin permutation count and the field-mismatch tax it bears. It also separates two costs that the framework keeps apart: the per-proof cost, which the precompile suite attacks, and the regeneration cost incurred when a relation changes, which only the static-circuit regime pays. Section~\ref{sec:evaluation} measures the per-proof term directly and characterises the regeneration term as a footprint: where the zkVM cost falls and how much of it the precompile removes, the resource frontier at which the unmitigated and the zero-knowledge-wrapped configurations cease to complete on the target hardware, and the recompile/setup time together with the enumerated regeneration footprint that separates the two regimes under a presentation, signature, or security-parameter change. The re-audit burden is reported as a component-count footprint and is not quantified as effort; the static baseline is Loquat/Aurora over $\mathbb{F}_{p_{127}}$ rather than same-scheme PLUM. The next chapter constructs the precompile suite that absorbs the dominant term this decomposition has located.
